\documentclass[11pt,a4paper]{article}

\usepackage[utf8]{inputenc}
\usepackage[english]{babel}
\usepackage{amsmath,amssymb,amsthm}
\usepackage{geometry}
\usepackage{hyperref}
\usepackage{bm}
\usepackage{physics}
\usepackage{mathtools}

\geometry{margin=2.5cm}

\title{Multi-level GRA Meta-Nullification in a Multiverse and All Regimes of GRA: Static, Cyclic, Chaotic, and Hybrid}
\author{AAA AAA}
\date{\today}

\begin{document}
\maketitle

\begin{abstract}
We formalize the multi-level architecture of GRA Meta-Nullification (Geometric Recursive Analytics) in terms of a multiverse of cognitive systems with an explicit hierarchy of goals $G_0,\dots,G_K$, a ``foam'' functional measuring cognitive inconsistency, and a global multiverse functional $J_{\text{multiverse}}$ whose minimization leads to a cognitive vacuum state. [web:8][web:20]
On top of this construction we distinguish four fundamental dynamical regimes of GRA --- static, cyclic, chaotic, and hybrid --- and give precise mathematical criteria for their applicability via linearization spectra and Lyapunov exponents. [web:19][web:20]
We illustrate the applicability of multi-level GRA by sketching a war-time multiverse of combat robots as a real-world example of a multi-scale control system. [web:18][web:13]
\end{abstract}

\section{Introduction}

GRA Nullification (Geometric Recursive Analytics) is a family of methods for minimizing ``foam'', understood as a functional of cognitive inconsistency between system states and a hierarchy of goals. [web:8][web:20]
Originally, GRA was formulated for a single abstraction level, but real-world cognition and control inevitably live in a multiverse of interconnected domains, from local physical processes to global political and economic structures. [web:8]
In this paper we:
(1) extend GRA to a multi-level multiverse architecture;
(2) formalize static, cyclic, chaotic, and hybrid GRA regimes;
(3) connect the formalism to a concrete example: a war-time multiverse of combat robots. [web:8][web:18]

\section{GRA Multiverse: Levels and Indexing}

\subsection{Hierarchy of meta-systems}

We consider an explicit hierarchy of meta-systems indexed by levels $l=0,\dots,K$:
\begin{itemize}
  \item Level $0$ (local): local nullifiers for individual domains (e.g., single robots or subsystems). [web:18]
  \item Level $1$ (meta): coordination of domains inside one meta-system (a platoon, a logistics cluster, a local sector). [web:18]
  \item Level $2$ (meta-meta): coordination of several meta-systems (an operational theater, a front line). [web:8]
  \item Level $K$ (multiversal): global coordination of the entire hierarchy (the full cognitive multiverse). [web:8]
\end{itemize}

Each object is identified by a multi-index
\[
\bm{a} = (a_0,a_1,\dots,a_k),
\]
where $a_0$ is the index of a local domain, $a_1$ is the index of a meta-system, and $a_k$ is the index at level $k$. [web:8]

\subsection{State spaces}

For level $l$ we define the tensor state space
\[
\mathcal{H}^{(l)} = \bigotimes_{\bm{a} : \dim(\bm{a})=l} \mathcal{H}^{(\bm{a})},
\]
where $\mathcal{H}^{(\bm{a})}$ is the state space of the local subsystem with index $\bm{a}$. [web:8]
The full multiverse state space is
\[
\mathcal{H}_{\text{multiverse}} = \bigotimes_{l=0}^K \mathcal{H}^{(l)}.
\]

A multiverse state is given by a family
\[
\bm{\Psi} = \{\Psi^{(\bm{a})}\}_{\bm{a}\in\mathcal{I}},
\]
where $\mathcal{I}$ is the set of all multi-indices. [web:8]

\section{Goal hierarchy and foam functional}

\subsection{Goals across levels}

Each level $l$ has an associated goal $G_l$:
\begin{itemize}
  \item Local goals: $G_0^{(\bm{a})}$ for each domain (e.g., successful mission of an individual robot). [web:18]
  \item Meta-goals: $G_1^{(\bm{b})}$ for coordinating domains inside a meta-system (logistics balance in a company / battalion / sector). [web:18]
  \item Meta-meta goals: $G_2^{(\bm{c})}$ for coordinating different meta-systems (resource balance across operational sectors). [web:8]
  \item Global goal: $G_K$ (strategic stability of the multiverse). [web:8]
\end{itemize}

\subsection{Foam at a single level}

Let $\Psi^{(l)}$ be the joint state of level $l$.
We define the foam functional at level $l$ as
\[
\Phi^{(l)}(\Psi^{(l)}, G_l) =
\sum_{\bm{a}\neq\bm{b} \atop \dim(\bm{a})=\dim(\bm{b})=l}
\big|\langle \Psi^{(\bm{a})} \vert \mathcal{P}_{G_l} \vert \Psi^{(\bm{b})} \rangle \big|^2,
\]
where $\mathcal{P}_{G_l}$ is the projector onto the solution space of goal $G_l$. [web:8]

Intuitively, $\Phi^{(l)}$ measures how strongly the subsystems of level $l$ disagree once projected onto the common goal $G_l$.

\subsection{Multi-level multiverse functional}

We define level-$l$ functionals recursively:
\begin{align}
J^{(0)}(\Psi^{(\bm{a})}) &= J_{\text{loc}}(\Psi^{(\bm{a})}; G_0^{(\bm{a})}), \\
J^{(l)}(\Psi^{(\bm{a})}) &=
\sum_{\bm{b}\prec\bm{a} \atop \dim(\bm{b})=l-1} J^{(l-1)}(\Psi^{(\bm{b})}) +
\Phi^{(l)}(\Psi^{(\bm{a})}, G_l^{(\bm{a})}),
\end{align}
where $\bm{b}\prec\bm{a}$ means that $\bm{b}$ is a subsystem of $\bm{a}$. [web:8]

The full super-meta multiverse functional is
\[
J_{\text{multiverse}}(\bm{\Psi}) =
\sum_{l=0}^K \Lambda_l \sum_{\dim(\bm{a})=l} J^{(l)}(\Psi^{(\bm{a})}),
\]
with level weights
\[
\Lambda_l = \lambda_0 \alpha^l,\quad 0<\alpha<1.
\]
This implements exponential decay of higher-level contributions and ensures a controlled influence of global constraints on local dynamics. [web:8]

Minimizing $J_{\text{multiverse}}$ corresponds to approaching a cognitive vacuum state, where interpretational artifacts are minimized across all levels. [web:20]

\section{Multiverse nullification theorem}

\subsection{Conditions}

Assume the following conditions:

\begin{enumerate}
  \item Full commutativity of projectors:
  \[
  [\mathcal{P}_{G_l^{(\bm{a})}}, \mathcal{P}_{G_m^{(\bm{b})}}] = 0,\quad \forall\,\bm{a},\bm{b},\, l,m.
  \]
  \item Hierarchical consistency of goals:
  \[
  \mathcal{P}_{G_l^{(\bm{a})}} =
  \bigotimes_{\bm{b}\prec\bm{a}} \mathcal{P}_{G_{l-1}^{(\bm{b})}},
  \quad l\geq 1.
  \]
  \item Space completeness:
  \[
  \dim(\mathcal{H}_{\text{multiverse}}) \geq \prod_{l=0}^K N_l,
  \]
  where $N_l$ is the number of subsystems at level $l$. [web:8]
\end{enumerate}

\subsection{Theorem and proof sketch}

\textbf{Theorem (multiverse nullification).}
If conditions (1)�(3) hold, then there exists a state $\bm{\Psi}^*$ such that
\[
\Phi^{(l)}(\Psi^{(l)*}, G_l) = 0,\quad \forall\, l = 0,\dots,K.
\]

\textit{Proof sketch.}
The base case $l=0$ follows from the local nullification theorem, i.e., there exist eigenvectors $\Psi^{(\bm{a})*}$ of the projectors $\mathcal{P}_{G_0^{(\bm{a})}}$ with zero foam. [web:8]
Induction step: assume foam is nullified at level $l-1$.
From hierarchical consistency we have
\[
\mathcal{P}_{G_l^{(\bm{a})}} = \bigotimes_{\bm{b}\prec\bm{a}} \mathcal{P}_{G_{l-1}^{(\bm{b})}}.
\]
Then the tensor product state
\[
\Psi^{(\bm{a})*} = \bigotimes_{\bm{b}\prec\bm{a}} \Psi^{(\bm{b})*}
\]
is an eigenvector of $\mathcal{P}_{G_l^{(\bm{a})}}$ with eigenvalue $1$. In this eigenbasis all off-diagonal elements vanish, thus yielding $\Phi^{(l)}=0$. [web:8]
Iterating this argument up to level $K$ produces global multiverse nullification.

\section{Algorithm GRA\_Multiverse\_Nullification}

\subsection{Recursive algorithm}

We describe a recursive algorithm that computes a minimal-foam state:

\begin{verbatim}
Algorithm GRA_Multiverse_Nullification
Input: Hierarchy of K+1 levels, goals {G_l}, parameters {?_l}
Output: Fully coherent multiverse state ?*

Function Nullify_Level(level l, state ?):

  if l = 0:
    return Local_Nullification(?, G_0)  # Base case

  else:
    subsystems = Decompose(?, l-1)

    for each sub in subsystems:
      sub = Nullify_Level(l-1, sub)

    ?_assembled = Recompose(subsystems)

    while ?^(l)(?_assembled, G_l) > ?_l:
      grad = ?_? ?^(l)(?_assembled, G_l)
      ?_assembled = ?_assembled - ?_l � grad

    return ?_assembled

Main procedure:

  ?_init = Initialize_Multiverse()
  ?_final = Nullify_Level(K, ?_init)
  return ?_final
\end{verbatim}

Parallel optimization is captured by
\[
\frac{\partial J_{\text{multiverse}}}{\partial \Psi^{(\bm{a})}} =
\Lambda_l \frac{\partial \Phi^{(l)}}{\partial \Psi^{(\bm{a})}}
+
\sum_{\bm{b}\succ\bm{a}} \Lambda_{l+1} \frac{\partial \Phi^{(l+1)}}{\partial \Psi^{(\bm{a})}},
\]
leading to the update rule
\[
\Psi^{(\bm{a})}(t+1) = \Psi^{(\bm{a})}(t) -
\eta \left[
\Lambda_l \frac{\partial \Phi^{(l)}}{\partial \Psi^{(\bm{a})}} +
\sum_{\bm{b}\succ\bm{a}} \Lambda_{l+1} \frac{\partial \Phi^{(l+1)}}{\partial \Psi^{(\bm{a})}}
\right]. [web:8]
\]

\section{Three fundamental GRA regimes}

\subsection{Static GRA (fixed point)}

Consider a dynamical system
\[
\dot{x} = f(x),\quad x\in\mathbb{R}^n,
\]
with an equilibrium $x_*$ such that $f(x_*)=0$.
Foam is given by a positive definite function, e.g.,
\[
\Phi(x) = \|x - x_{\text{goal}}\|^2.
\]
Static GRA aims for $x(t)\to x_*$ with $\Phi$ minimized.
Applicability criterion: all eigenvalues of the Jacobian
\[
A = \left.\frac{\partial f}{\partial x}\right|_{x_*}
\]
satisfy
\[
\Re(\lambda_i(A)) < 0,
\]
which ensures asymptotic stability. [web:20]

In the simplest case evolution follows gradient descent
\[
\dot{x} = -\nabla_x \Phi(x) + \text{noise}.
\]

\subsection{Cyclic GRA (limit cycle)}

Cyclic GRA has a limit cycle
\[
\Gamma = \{x(t): x(t+T)=x(t)\}
\]
as its target attractor. Foam is defined as the mean squared deviation from a reference cycle $x_{\text{ref}}(t)$:
\[
\Phi_{\text{cycle}} = \frac{1}{T}\int_0^T \|x(t)-x_{\text{ref}}(t)\|^2 dt.
\]
The canonical mathematical criterion is a Hopf bifurcation: the linearized system at an equilibrium has a pair of purely imaginary eigenvalues $\pm i\omega$, all others with $\Re<0$, and the first Lyapunov coefficient $l_1<0$, guaranteeing the birth of a stable limit cycle. [web:20]

\subsection{Chaotic GRA (strange attractor)}

In the chaotic regime the system exhibits deterministic chaos: trajectories are aperiodic, sensitive to initial conditions, yet confined to a strange attractor. [web:19]
Foam is defined as a statistical quantity, such as
\[
\Phi_{\text{chaos}} = \mathbb{E}\,\|x - x_{\text{ref}}\|^2,
\]
or via entropy-like functionals.
A standard criterion for chaos is a positive largest Lyapunov exponent $\lambda_1>0$. [web:19]

Chaos control can be implemented through feedback of the form
\[
\dot{x} = f(x) + K\big(x(t-\tau) - x(t)\big),
\]
as in Pyragas control, stabilizing unstable periodic orbits embedded in the chaotic attractor. [web:19]

\section{Hybrid GRA in a multi-level hierarchy}

\subsection{Mixing regimes across levels}

In a multi-level architecture it is natural to assign different GRA regimes to different $l$:
\begin{itemize}
  \item Lower (fast) levels: static GRA.
  \item Intermediate levels: cyclic GRA.
  \item Top levels: chaotic or statistical GRA.
\end{itemize}

The hybrid functional is
\[
J_{\text{hybrid}} = \sum_{l=0}^K \Lambda_l\,\Phi_{\text{type}(l)}^{(l)},
\]
where $\Phi_{\text{type}(l)}^{(l)}$ is computed according to the regime chosen for level $l$. [web:20]

\section{Example: a war-time multiverse of combat robots}

\subsection{Level indexing}

Consider a war-time multiverse in which:
\begin{itemize}
  \item Level $0$: individual robots or micro-swarms (logistics, medevac, strike). [web:18]
  \item Level $1$: tactical units / local sectors (clusters around logistics hubs and kill zones). [web:18]
  \item Level $2$: operational theaters (front directions). [web:8]
  \item Level $3$: global war strategy (casualties, industry, politics). [web:8]
\end{itemize}
A multi-index $\bm{a}=(a_0,a_1,a_2,a_3)$ encodes a robot, its unit, its theater, and the strategic context.

\subsection{Goals and regimes}

\begin{itemize}
  \item $G_0^{(\bm{a})}$: successful mission of the robot with minimized kill probability and strict deadlines; static GRA is appropriate here. [web:18]
  \item $G_1^{(\bm{b})}$: maintaining a prescribed fraction of logistics and evacuation by robots for the local unit, with minimal route conflicts; cyclic GRA stabilizes repeating mission cycles. [web:18]
  \item $G_2^{(\bm{c})}$: resource balance between sectors and robustness of supply chains and operational tempo. [web:8]
  \item $G_3$: minimizing human casualties for a given probability of strategic success, while keeping industry and political multiverse within viable bounds; chaotic GRA manages global chaos statistically. [web:8]
\end{itemize}

The foam at each level penalizes inconsistency of local plans with respect to the level goal and inconsistency between levels through the recursive structure of $J^{(l)}$. [web:8][web:20]

\subsection{Hybrid architecture}

In this war-time multiverse, hybrid GRA implements:
\begin{itemize}
  \item static stabilization of individual robots and low-level subsystems,
  \item cyclic stabilization of logistics and operational cycles at intermediate levels,
  \item chaotic control of strategic processes via statistical risk minimization at the top. [web:18][web:20]
\end{itemize}

Thus, the GRA Meta-Nullification multiverse turns a complex, conflict-driven environment into a principled optimization problem under an explicit hierarchy of ethical and strategic constraints. [web:8]

\section{Conclusion}

Multi-level GRA Meta-Nullification in a multiverse scales the nullification principle to arbitrary hierarchical structures, from local dynamical systems to global cognitive and political multiverses. [web:8][web:20]
The classification of GRA regimes via linearization spectra and Lyapunov exponents provides sharp applicability criteria for static, cyclic, chaotic, and hybrid regimes. [web:19][web:20]
The combat robot multiverse example shows how this formalism can be used to design real-world multi-level control systems under extreme uncertainty and adversarial pressure. [web:18]

\end{document}